When you create an interface or gateway endpoint, you can attach an
endpoint policy to it that controls access to the service to which you
are connecting. A VPC endpoint policy is an IAM resource policy that you
attach to an endpoint. If you do not attach a policy when you create an
endpoint, AWS attaches a default policy for you that allows full access
to the service through the endpoint.

You can create a policy that restricts access to specific S3 buckets
only. This is useful if you only want certain S3 Buckets to be
accessible through the endpoint.

In this section you will create a VPC endpoint policy that restricts
access to the S3 bucket specified in the VPC endpoint policy.

\begin{figure}
\centering
\includegraphics{5-Workshop/5.5-Policy/s3-bucket-policy.png}
\caption{endpoint diagram}
\end{figure}

\vspace{0.5em}\noindent\textbf{Connect to an EC2 instance and verify connectivity to
S3}\par

\begin{enumerate}
\def\labelenumi{\arabic{enumi}.}
\tightlist
\item
  Start a new AWS Session Manager session on the instance named
  Test-Gateway-Endpoint. From the session, verify that you can list the
  contents of the bucket you created in Part 1: Access S3 from VPC:
\end{enumerate}

\begin{verbatim}
aws s3 ls s3://\<your-bucket-name\>
\end{verbatim}

The bucket contents include the two 1 GB files uploaded in earlier.

\begin{enumerate}
\def\labelenumi{\arabic{enumi}.}
\setcounter{enumi}{1}
\tightlist
\item
  Create a new S3 bucket; follow the naming pattern you used in Part 1,
  but add a `-2' to the name. Leave other fields as default and click
  create
\end{enumerate}

Successfully create bucket

\begin{enumerate}
\def\labelenumi{\arabic{enumi}.}
\setcounter{enumi}{2}
\tightlist
\item
  Navigate to: Services \textgreater{} VPC \textgreater{} Endpoints,
  then select the Gateway VPC endpoint you created earlier. Click the
  Policy tab. Click Edit policy.
\end{enumerate}

The default policy allows access to all S3 Buckets through the VPC
endpoint.

\begin{enumerate}
\def\labelenumi{\arabic{enumi}.}
\setcounter{enumi}{3}
\tightlist
\item
  In Edit Policy console, copy \& Paste the following policy, then
  replace yourbucketname-2 with your 2nd bucket name. This policy will
  allow access through the VPC endpoint to your new bucket, but not any
  other bucket in Amazon S3. Click Save to apply the policy.
\end{enumerate}

\emph{Detailed configuration is omitted from the condensed PDF; see the
web report for the complete listing.}

Successfully customize policy

\begin{enumerate}
\def\labelenumi{\arabic{enumi}.}
\setcounter{enumi}{4}
\tightlist
\item
  From your session on the Test-Gateway-Endpoint instance, test access
  to the S3 bucket you created in Part 1: Access S3 from VPC
\end{enumerate}

\begin{verbatim}
aws s3 ls s3://<yourbucketname>
\end{verbatim}

This command will return an error because access to this bucket is not
permitted by your new VPC endpoint policy:

\begin{enumerate}
\def\labelenumi{\arabic{enumi}.}
\setcounter{enumi}{5}
\tightlist
\item
  Return to your home directory on your EC2 instance
  \texttt{cd\textasciitilde{}}
\end{enumerate}

\begin{itemize}
\tightlist
\item
  Create a file \texttt{fallocate\ -l\ 1G\ test-\allowbreak{}bucket2.\allowbreak{}xyz}
\item
  Copy file to 2nd bucket
  \texttt{aws\ s3\ cp\ test-bucket2.xyz\ s3://\textless{}your-2nd-bucket-name\textgreater{}}
\end{itemize}

This operation succeeds because it is permitted by the VPC endpoint
policy.

\begin{itemize}
\tightlist
\item
  Then we test access to the first bucket by copy the file to 1st bucket
  \texttt{aws\ s3\ cp\ test-bucket2.xyz\ s3://\textless{}your-1st-bucket-name\textgreater{}}
\end{itemize}

This command will return an error because access to this bucket is not
permitted by your new VPC endpoint policy.

\vspace{0.5em}\noindent\textbf{Part 3 Summary:}\par

In this section, you created a VPC endpoint policy for Amazon S3, and
used the AWS CLI to test the policy. AWS CLI actions targeted to your
original S3 bucket failed because you applied a policy that only allowed
access to the second bucket you created. AWS CLI actions targeted for
your second bucket succeeded because the policy allowed them. These
policies can be useful in situations where you need to control access to
resources through VPC endpoints.
